\subsection{The connected plaquette correlation on a finite torus}
\label{ssec:rEtwo}

This subsection records what the Wilson measure of a finite torus gives for
the connected two-point function of the plaquette field $\Phi$ of
Sec.~\ref{ssec:limitE} at every coupling, and what it gives in the
strong-coupling range. The results are unconditional. They do not prove
either half of the two-sided law (AF) of Remark~\ref{rem:limitEafstatus};
the last part restates the lower half of (AF) along an axis in spectral
terms and locates the step at which an argument for it leaves the range
of the tools used here.

\emph{Conventions.} Throughout, $G\subset\U(N)$ is a compact group with a
faithful representation $\rho$, by which it is identified with a closed
subgroup of $\U(N)$, so that $G$ is a compact Lie group; $\chi=\Tr\rho$ is
its character, $\bar\chi$ the character of the dual representation, and
$\delta=1$ if $\chi$ is real and $\delta=0$ otherwise. From
Lemma~\ref{rEtwo:polymer} on, $\rho$ is irreducible, as it is for the
defining representation of every compact simple group. The Wilson measure
$\mu_L$ of $\Lambda_L=(\Z/L\Z)^4$, $L\ge6$ even, has density
$\exp(N\beta\sum_p\ReTr U_p)$ with respect to the product Haar measure on
the links, $\beta>0$, as in Sec.~\ref{sec:intro}, and $\langle\cdot
\rangle_L$ is its expectation. For a plaquette $p$ put $F_p=N-\ReTr U_p
\ge0$, and for a site $x$ let $\Phi_x=\sum_{p\ni x}F_p$ be the sum over the
$24$ plaquettes with a corner at $x$; this is the field $\Phi$ of
Sec.~\ref{ssec:limitE} with $r=1$. Put
\[
c_L=\langle\Phi_0\rangle_L,\qquad G_L(z)=\langle\Phi_0\Phi_z\rangle_L-c_L^2
\qquad(z\in\Lambda_L).
\]
By the invariance of $\mu_L$ under translations and the hypercubic group
(proof of Lemma~\ref{rE:torus}(ii)), $G_L(z)=\langle\Phi_x\Phi_{x+z}
\rangle_L-c_L^2$ for every $x$, $G_L(gz)=G_L(z)$ for $g\in B_4$, and
$G_L(-z)=G_L(z)$. The torus limits $\omega$ of Lemma~\ref{rE:torus} are
weak limits of $\mu_L$ along subsequences; for them $c=\omega(\Phi_0)$,
$G(z)=\omega(\Phi_0\Phi_z)-c^2$, and $G(z)=\lim G_{L}(z)$ along the
subsequence, since $\Phi_0\Phi_z$ is a continuous function of finitely
many links.

\subsubsection{The transfer matrix of the torus}

For $\beta>0$ and an irreducible representation $R$ of $G$ let
\[
c_R(\beta)=\int_Ge^{N\beta\ReTr g}\,\overline{\chi_R(g)}\,dg,\quad
e^{N\beta\ReTr g}=\sum_Rc_R(\beta)\chi_R(g),
\]
the series converging absolutely and uniformly, this being
the Peter--Weyl expansion of a smooth class function on a compact Lie
group. Expanding the exponential, $c_R(\beta)=\sum_{n\ge0}(N\beta/2)^n
m^{(n)}_R/n!$, where $m^{(n)}_R\in\Z_{\ge0}$ is the multiplicity of $R$ in
$(\rho\oplus\bar\rho)^{\otimes n}$; every $R$ occurs in some tensor power
because $\rho$ is faithful (Stone--Weierstrass and Schur orthogonality, as
in the proof of Proposition~\ref{prop:flowpositive}), so
\begin{equation}\label{rEtwo:eqcpos}
c_R(\beta)>0\qquad\text{for every $R$ and every }\beta>0 .
\end{equation}

Fix the time direction $0$. Let $\Lambda^3=(\Z/L\Z)^3$ be the spatial
torus, $E^3$ its links, $P^3$ its plaquettes, $\mathfrak h=L^2(G^{E^3})$
with respect to the product Haar measure, and let the gauge group
$\mathcal G=G^{\Lambda^3}$ act on $G^{E^3}$ by $(g\cdot U)_{xy}=g_xU_{xy}
g_y^{-1}$ and on $\mathfrak h$ by the unitaries $(R_g\psi)(U)=\psi(g\cdot
U)$. Let $P=\int_{\mathcal G}R_g\,dg$ be the orthogonal projection onto
the gauge-invariant functions and $\mathfrak h_P=P\mathfrak h$. Define
the kernel
\begin{align*}
T(U,U')&=D(U)\,K(U,U')\,D(U'),\\
D(U)&=\exp\Bigl(\frac{N\beta}2\sum_{p\in P^3}\ReTr U_p\Bigr),\\
K(U,U')&=\prod_{l\in E^3}\exp\bigl(N\beta\ReTr(U_lU_l'^{-1})\bigr),
\end{align*}
and the operator $(T\psi)(U)=\int T(U,U')\psi(U')\,dU'$ on $\mathfrak h$.
This is the transfer matrix of L\"uscher~\cite{Lus77} in the temporal
gauge, and of Osterwalder and Seiler~\cite{OSei78}; its positivity is
recalled in the proof of Proposition~\ref{prop:flowpositive}. What is
needed here beyond that proof is the trace formula on a torus of finite
time extent and the form of the plaquette insertions.

\begin{lemma}[the transfer matrix]\label{rEtwo:transfer}
\textup{(i)} $T$ is a self-adjoint Hilbert--Schmidt operator on
$\mathfrak h$ with $\ip\psi{T\psi}\ge0$ for all $\psi$, and $T$ is
injective. It commutes with every $R_g$, with $P$, and with the unitaries
induced by the translations and the cubic symmetries of $\Lambda^3$.
Consequently $T$ maps $\mathfrak h_P$ into itself, and $\mathfrak h_P$ has
an orthonormal basis $(e_j)_{j\ge0}$ of eigenvectors of $T$ with
eigenvalues $\lambda_j>0$, $\lambda_j\to0$.

\textup{(ii)} $Z_L=\int\exp(N\beta\sum_p\ReTr U_p)\,dU=\Tr(PT^L)$.

\textup{(iii)} Let $B$ be the operator on $\mathfrak h$ with kernel
\begin{align*}
B(U,U')&=\int\Phi(U,V,U')\,T(U,V)\,T(V,U')\,dV,\\
\Phi(U,V,U')&=\sum_{p\in P^3,\,p\ni0}(N-\ReTr V_p)\\
&\quad+\sum_{l\in E^3,\,l\ni0}\bigl[(N-\ReTr(U_lV_l^{-1}))\\
&\qquad\qquad+(N-\ReTr(V_lU_l'^{-1}))\bigr].
\end{align*}
Then $B$ is self-adjoint and Hilbert--Schmidt, commutes with every
$R_g$ and with $P$, and for $2\le t\le L-2$
\[
c_L=\frac{\Tr(PT^{L-2}B)}{Z_L},\quad\langle\Phi_0\Phi_{te_0}\rangle_L=
\frac{\Tr\bigl(PT^{L-t-2}BT^{t-2}B\bigr)}{Z_L}.
\]
\end{lemma}

\begin{proof}
(i) The kernel is continuous and real, and symmetric because $K(U,U')=
K(U',U)$, which follows from $\ReTr(U'U^{-1})=\ReTr((UU'^{-1})^{-1})=
\ReTr(UU'^{-1})$ for unitary matrices; so $T$ is self-adjoint and
Hilbert--Schmidt. Write $T=\mathsf DK\mathsf D$ with $\mathsf D$ the
multiplication by the bounded function $D$, which has values in a compact
subinterval of $(0,\infty)$, so that $\mathsf D$ is invertible. The
operator $K$ is the tensor product over $l\in E^3$ of the convolution
operators $K_l$ on $L^2(G)$ with kernel $k_\beta(U_lU_l'^{-1})$,
$k_\beta(g)=e^{N\beta\ReTr g}=\sum_Rc_R(\beta)\chi_R(g)$. After the
substitution $V=W^{-1}U$,
\begin{align*}
\int k_\beta(UV^{-1})R(V)_{ij}\,dV&=\sum_m\Bigl(\int k_\beta(W)\,
\overline{R(W)_{mi}}\,dW\Bigr)R(U)_{mj}\\
&=\frac{c_R(\beta)}{d_R}\,R(U)_{ij}
\end{align*}
by the Schur orthogonality relations, so $K_l$ acts on the matrix
coefficients of each irreducible $R$, which form an orthogonal basis of
$L^2(G)$, as multiplication by $c_R(\beta)/d_R>0$ \eqref{rEtwo:eqcpos}.
Hence $K$ is positive and injective, its eigenvalues on the product basis
being the products $\prod_lc_{R_l}(\beta)/d_{R_l}>0$, and $\ip\psi{T\psi}=
\ip{\mathsf D\psi}{K\mathsf D\psi}\ge0$; $T\psi=0$ gives $K\mathsf D\psi=0$,
so $\psi=0$. The kernel satisfies $T(g\cdot U,g\cdot U')=T(U,U')$ because
every $\ReTr$ in it is the trace of a product around a closed loop of
links, so $T$ commutes with $R_g$, hence with $P$; the same holds for the
lattice symmetries, which permute links and plaquettes and preserve
$\ReTr U_p$ under reversal of orientation. A positive compact injective
operator on the separable Hilbert space $\mathfrak h_P$ has an orthonormal
eigenbasis with positive eigenvalues tending to $0$.

(ii) Write the links of $\Lambda_L$ as spatial links $U_{(x,s),i}$ and
temporal links $U_{(x,s),0}$, $s\in\Z/L\Z$. Given a configuration, define
$h_{(x,0)}=1$ and $h_{(x,s)}=U_{(x,0),0}U_{(x,1),0}\cdots U_{(x,s-1),0}$
for $1\le s\le L-1$, and let $U'=h\cdot U$. Then $U'_{(x,s),0}=1$ for
$s\le L-2$ and $U'_{(x,L-1),0}=g_x:=\prod_{s=0}^{L-1}U_{(x,s),0}$ (ordered
product). The map $U\mapsto(U'_{(x,s),i},\,U_{(x,s),0}\ (s\le L-2),\,g_x)$
preserves the product Haar measure: at each site the map $(U_0,\dots,U_{
L-1})\mapsto(U_0,\dots,U_{L-2},U_0\cdots U_{L-1})$ of $G^L$ does, and for
fixed temporal links the spatial links are transformed by a gauge
transformation, which does. The Wilson density and every gauge-invariant
observable depend only on $(U'_{(x,s),i},g)$; the variables $U_{(x,s),0}$,
$s\le L-2$, integrate to $1$. With $U_s$ the spatial configuration of
slice $s$, the temporal plaquette at the link $l=\{x,x+e_i\}$ between the
slices $s$ and $s+1\le L-1$ is $U_{s,l}U_{s+1,l}^{-1}$, and between the
slices $L-1$ and $0$ it is $U_{L-1,l}g_{x+e_i}U_{0,l}^{-1}g_x^{-1}=U_{L-1,l}
((g\cdot U_0)_l)^{-1}$. Distributing each spatial factor $D(U_s)^2$ over
the two adjacent temporal steps gives
\begin{align*}
Z_L=\int_{\mathcal G}dg\int T(U_0,U_1)&\cdots T(U_{L-2},U_{L-1})\\
&\times T(U_{L-1},g\cdot U_0)\,\prod_sdU_s .
\end{align*}
The inner integral is $\int T^L(U_0,g\cdot U_0)\,dU_0=\Tr(T^LR_g)$, since
$T^L$ has a continuous symmetric kernel and is trace class ($L\ge2$), and
$\int\Tr(T^LR_g)\,dg=\Tr(T^LP)$.

(iii) Place $\Phi_0$ at the vertex $(0,1)$, so that it is a function of
the slices $0,1,2$: its $12$ spatial plaquettes lie in slice $1$, and its
$12$ temporal plaquettes are $U_{0,l}U_{1,l}^{-1}$ and $U_{1,l}U_{2,l}^{-1}$
for the six spatial links $l$ at $0$; this is $\Phi(U_0,U_1,U_2)$. Inserting
it into the integral of (ii), with the step carrying $g$ chosen among the
$L-2\ge1$ steps not touched by the observable (the relabelling of the
times is a symmetry of $\mu_L$), gives $\langle\Phi_0\rangle_LZ_L=
\Tr(PT^{L-2}B)$. For the two-point function place the second copy at
$(0,t+1)$, a function of the slices $t,t+1,t+2$; for $2\le t\le L-2$ the
two copies are separated by $t-2\ge0$ and $L-t-2\ge0$ transfer steps,
$(t-2)+(L-t-2)=L-4\ge2$, and the step carrying $g$ is chosen in one of the
two gaps. The kernel $B$ is real and symmetric because $\Phi(U,V,U')=
\Phi(U',V,U)$ and $T$ is symmetric, and gauge covariant because $\Phi$ and
$T$ are.
\end{proof}

\subsubsection{Positivity, convexity and the dichotomy along an axis}

Let $(\lambda_j,e_j)$ be as in Lemma~\ref{rEtwo:transfer}(i), $B_{jk}=
\ip{e_j}{Be_k}$, $\pi_j=\lambda_j^L/Z_L$ and $b_j=B_{jj}/\lambda_j^2$. Then
$\sum_j\pi_j=1$, $\pi_j>0$, and $Z_Lc_L=\sum_j\lambda_j^{L-2}B_{jj}$ gives
$c_L=\sum_j\pi_jb_j$. Put
\[
V_L=\sum_j\pi_jb_j^2-\Bigl(\sum_j\pi_jb_j\Bigr)^2\ge0 .
\]

\begin{theorem}[the axial two-point function on a torus]\label{rEtwo:axial}
Let $\beta>0$ and $L\ge6$ even. For $2\le t\le L-2$ and every direction
$\mu$,
\[
G_L(te_\mu)=G_L(te_0)=V_L+\frac1{Z_L}\sum_{j\ne k}\lambda_j^{L-t-2}
\lambda_k^{t-2}\,|B_{jk}|^2 ,
\]
the series converging absolutely. Consequently:

\textup{(i)} $G_L(te_0)\ge V_L\ge0$ and $G_L(te_0)=G_L((L-t)e_0)$;

\textup{(ii)} $t\mapsto G_L(te_0)$ is log-convex on $\{2,\dots,L-2\}$, that
is $G_L(te_0)^2\le G_L((t-1)e_0)\,G_L((t+1)e_0)$ for $3\le t\le L-3$, and
nonincreasing on $\{2,\dots,L/2\}$;

\textup{(iii)} either $G_L(te_0)>0$ for every $t\in\{2,\dots,L-2\}$, or
$G_L(te_0)=0$ for every such $t$; the second alternative holds if and
only if $B=c_LT^2$ on $\mathfrak h_P$.

\textup{(iv)} In the second alternative of \textup{(iii)}, $G_L(z)=0$ for
every $z\in\Lambda_L$ having a coordinate $z_\mu$ with $2\le|z_\mu|\le
L-2$, and $\langle\Phi_0\rangle$ on $(\Z/L'\Z)\times\Lambda^3$ equals $c_L$
for every even $L'\ge6$.
\end{theorem}

\begin{proof}
The equality $G_L(te_\mu)=G_L(te_0)$ is the $B_4$-invariance of $\mu_L$
and $\Phi$. By Lemma~\ref{rEtwo:transfer}(iii) and the eigenbasis of
$\mathfrak h_P$, using that $B$ commutes with $P$ so that $PT^{L-t-2}BT^{
t-2}B=PT^{L-t-2}(PBP)T^{t-2}(PBP)$,
\begin{align*}
Z_L\langle\Phi_0\Phi_{te_0}\rangle_L&=\sum_{j,k}\lambda_j^{L-t-2}\lambda_k^{
t-2}B_{jk}B_{kj}\\
&=\sum_{j,k}\lambda_j^{L-t-2}\lambda_k^{t-2}|B_{jk}|^2 ,
\end{align*}
since $B$ is self-adjoint; the terms are nonnegative and their sum is the
trace of $PT^{L-t-2}BT^{t-2}B$, which is trace class: $T^aB$ is
Hilbert--Schmidt for $a\ge1$ and the product of two Hilbert--Schmidt
operators is trace class, and if $t-2=0$ or $L-t-2=0$ the other exponent is
$L-4\ge2$ and $BT^{L-4}B=(BT^{L/2-2})(T^{L/2-2}B)$. The diagonal terms
give $\sum_j\lambda_j^{L-4}B_{jj}^2/Z_L=\sum_j\pi_jb_j^2$, and subtracting
$c_L^2=(\sum_j\pi_jb_j)^2$ yields the formula.

(i) Both parts are nonnegative, and $V_L\ge0$ is the variance of $b$ under
$\pi$. Exchanging $j$ and $k$ in the sum shows the symmetry $t\mapsto L-t$.

(ii) Each term $\lambda_j^{L-t-2}\lambda_k^{t-2}|B_{jk}|^2=C_{jk}(\lambda_k/
\lambda_j)^t$ with $C_{jk}\ge0$ is log-convex in $t$, a constant is
log-convex, and sums of log-convex functions are log-convex (by the
Cauchy--Schwarz inequality for the measure with weights $C_{jk}$).
Grouping $(j,k)$ with $(k,j)$ gives $(\lambda_j\lambda_k)^{L/2-2}\bigl(x^{
L/2-t}+x^{t-L/2}\bigr)|B_{jk}|^2$ with $x=\lambda_j/\lambda_k$, which is
nonincreasing in $t$ on $[2,L/2]$.

(iii) If $G_L(te_0)=0$ for one $t$, then $V_L=0$ and every term of the
off-diagonal sum vanishes; as all $\lambda_j>0$, $B_{jk}=0$ for all $j\ne k$,
so $Be_k=B_{kk}e_k=b_k\lambda_k^2e_k$, and $V_L=0$ with $\pi_k>0$ for all
$k$ forces $b_k=\sum_j\pi_jb_j=c_L$ for all $k$. Hence $B=c_LT^2$ on
$\mathfrak h_P$, and then the formula gives $G_L(te_0)=0$ for every $t$.
Conversely $B=c_LT^2$ on $\mathfrak h_P$ gives $B_{jk}=c_L\lambda_k^2
\delta_{jk}$, $b_k=c_L$, $V_L=0$.

(iv) Let $z=(t,\vec z)$ with $2\le t\le L-2$ and let $\tau_{\vec z}$ be
the unitary of the spatial translation by $\vec z$; it commutes with $T$
and $P$. The proof of Lemma~\ref{rEtwo:transfer}(iii) gives $Z_L\langle
\Phi_0\Phi_z\rangle_L=\Tr(PT^{L-t-2}BT^{t-2}\tau_{\vec z}B\tau_{\vec z}^{-1})$,
which under $B=c_LT^2$ equals $c_L^2\Tr(PT^L)=c_L^2Z_L$. The other
directions follow from the $B_4$-invariance. For the time extent $L'$,
Lemma~\ref{rEtwo:transfer}(ii),(iii) with $L'$ in place of $L$ give
$\langle\Phi_0\rangle=\Tr(PT^{L'-2}B)/\Tr(PT^{L'})=c_L$.
\end{proof}

\begin{remark}\label{rEtwo:axialremark}
Reflection positivity of $\mu_L$ in the plane $\{x_0=0\}$ through sites
(with its antipode) and in the plane $\{x_0=\frac12\}$ bisecting
links~\cite{OSei78} gives $G_L(te_0)\ge0$ for all $2\le t\le L-2$: for even
$t=2s$ the function $\Phi_{se_0}-c_L$ depends on links at times in
$[s-1,s+1]\subset[0,L/2]$ and is reflected onto $\Phi_{-se_0}-c_L$, and for
odd $t=2s+1$ the function $\Phi_{(s+1)e_0}-c_L$ depends on links at times
in $[s,s+2]\subset[1,L/2]$ and is reflected onto $\Phi_{-se_0}-c_L$.
Theorem~\ref{rEtwo:axial} gives in addition the floor $V_L$, the
log-convexity, the monotonicity and the exact form of the degenerate
alternative. The quantity $V_L$ is the variance, under the Gibbs weights
$\pi_j$ of the torus of time extent $L$, of the diagonal values $b_j$; it
does not depend on $t$.
\end{remark}

The torus limits inherit the structure. Let $\omega$ be a torus limit as
in Lemma~\ref{rE:torus}, and let $\vartheta'$ be the reflection in the
plane $\{x_0=\tfrac12\}$ bisecting links, for which $\omega$ is reflection
positive by Lemma~\ref{rE:torus}(iii). The reconstruction of
Appendix~\ref{app:transfer} applies to $\vartheta'$ as to $\vartheta$: the
positive-time functions are those of the links with both endpoints at
times $\ge1$, the time translation $\tau_1$ maps them into themselves, and
$\vartheta'\tau_t=\tau_{-t}\vartheta'$; it gives $\mathbb T'$, $\Omega'$ and
$F\mapsto\Psi'_F$ with $\ip{\Psi'_F}{\mathbb T'^n\Psi'_H}=\omega(\vartheta'
(\bar F)\tau_nH)$ for $n\ge0$. Let $\Psi'=\Psi'_{\Phi_{2e_0}-c}$ be the
vector of the centred field at the site $2e_0$, whose links have both
endpoints at times in $[1,3]$, and let $\sigma$ be the spectral measure of
$\mathbb T'$ at $\Psi'$, a finite measure on $[0,1]$.

\begin{corollary}[torus limits]\label{rEtwo:limit}
For every torus limit $\omega$ at $\beta\ge0$:

\textup{(i)} $G(te_0)\ge0$ for all $t\ge2$, $t\mapsto G(te_0)$ is
log-convex and nonincreasing on $\{2,3,\dots\}$, and for $n\ge0$
\[
G((3+n)e_0)=\ip{\Psi'}{\mathbb T'^n\Psi'}=\int_{[0,1]}\lambda^n\,d\sigma(\lambda);
\]

\textup{(ii)} either $G(te_0)>0$ for all $t\ge3$, or $G(te_0)=0$ for all
$t\ge4$; the second alternative holds if and only if $\sigma$ is
concentrated at $\lambda=0$;

\textup{(iii)} in the first alternative the local rates
\[
m(t)=\log\bigl(G(te_0)/G((t+1)e_0)\bigr),\qquad t\ge2,
\]
are nonnegative, nonincreasing in $t$, and converge to
$m_\Phi=-\log\max\operatorname{supp}\sigma\in[0,\infty)$, which satisfies
$m_\Phi\ge\log(1/\gamma')$ with $\gamma'=\|\mathbb T'|_{\Omega'^\perp}\|$,
the analogue of $\gamma(\omega)$ for the reflection $\vartheta'$.
\end{corollary}

\begin{proof}
(i) For fixed $t$, $G(te_0)=\lim G_L(te_0)$ along the subsequence, and
positivity, log-convexity on $\{2,\dots,L-2\}$ and monotonicity on
$\{2,\dots,L/2\}$ pass to the limit, the ranges exhausting $\{2,3,\dots\}$.
The identity for $n\ge0$ is the reconstruction: $\vartheta'(\Phi_{2e_0}-c)
=\Phi_{-e_0}-c$, $\tau_n(\Phi_{2e_0}-c)=\Phi_{(2+n)e_0}-c$, so $\omega(
\vartheta'(\bar F)\tau_nF)=G((3+n)e_0)$ with $F=\Phi_{2e_0}-c$, and this
equals $\ip{\Psi'_F}{\mathbb T'^n\Psi'_F}$.

(ii) If $G(t_0e_0)=0$ for some $t_0\ge4$, then $\int\lambda^{t_0-3}d\sigma
=0$, so $\sigma((0,1])=0$ and $G(te_0)=\int\lambda^{t-3}d\sigma=0$ for all
$t\ge4$; if $G(3e_0)=0$ then $\sigma=0$. If $\sigma((0,1])>0$ then
$\int\lambda^nd\sigma>0$ for all $n\ge0$.

(iii) The ratios $\int\lambda^{n+1}d\sigma/\int\lambda^nd\sigma$ of the
moments of a finite measure on $[0,1]$ with $\sigma((0,1])>0$ are at most
$1$, nondecreasing in $n$ by the Cauchy--Schwarz inequality, and converge
to $\max\operatorname{supp}\sigma$; so $m(t)$ is nonnegative, nonincreasing
and converges to $m_\Phi$ for $t\ge3$. For $t=2$ use the log-convexity and
monotonicity of (i) directly. The vector $\Psi'$ is orthogonal to
$\Omega'$ since $\omega(F)=0$, so $\operatorname{supp}\sigma\subset[0,\|
\mathbb T'|_{\Omega'^\perp}\|]$.
\end{proof}

\subsubsection{Strict positivity and the exceptional couplings}

The dichotomy of Theorem~\ref{rEtwo:axial}(iii) is decided for a single
spatial plaquette at every coupling.

\begin{theorem}[one plaquette]\label{rEtwo:plaquette}
Let $\beta>0$, $L\ge6$ even, let $p$ be a plaquette in a plane not
containing the direction $0$, and let $1\le t\le L-1$. Then
$\langle F_pF_{p+te_0}\rangle_L-\langle F_p\rangle_L^2>0$. The same holds
for every direction $\mu$ not in the plane of $p$ in place of $e_0$.
\end{theorem}

\begin{proof}
Place $p$ in the slice $0$ and let $M_p$ be the multiplication by $F_p$
on $\mathfrak h$, a bounded self-adjoint operator commuting with every
$R_g$. The proof of Lemma~\ref{rEtwo:transfer}(iii), with $M_p$ at the
slices $0$ and $t$ and the step carrying $g$ in one of the two gaps of
lengths $t\ge1$ and $L-t\ge1$, gives $Z_L\langle F_pF_{p+te_0}\rangle_L=
\Tr(PT^{L-t}M_pT^tM_p)=\sum_{j,k}\lambda_j^{L-t}\lambda_k^t|M_{jk}|^2$ with
$M_{jk}=\ip{e_j}{M_pe_k}$, and $Z_L\langle F_p\rangle_L=\sum_j\lambda_j^L
M_{jj}$. As in the proof of Theorem~\ref{rEtwo:axial}, the covariance is
the sum of the variance of $(M_{jj})_j$ under $(\pi_j)_j$ and of
$Z_L^{-1}\sum_{j\ne k}\lambda_j^{L-t}\lambda_k^t|M_{jk}|^2$, both
nonnegative, and it vanishes only if $M_{jk}=0$ for $j\ne k$ and $M_{jj}$
is constant, that is, only if $M_p\psi=c\psi$ for every $\psi\in\mathfrak
h_P$. The constant function $1$ lies in $\mathfrak h_P$, so $F_p$ would be
almost everywhere equal to $c$. The map $(U_l)_{l\in p}\mapsto U_p$
carries the product Haar measure on $G^4$ to the Haar measure of $G$,
which has full support, and $g\mapsto\ReTr g$ is continuous; so $\ReTr g=c$
would hold on a dense set, hence everywhere, in particular $c=\ReTr1=N$,
whereas $\ReTr g<N$ for $g\ne1$ and $G\ne\{1\}$. The statement for $\mu$
follows by a cubic symmetry.
\end{proof}

For the field $\Phi$ the first ingredient is analyticity in the coupling.

\begin{lemma}\label{rEtwo:analytic}
For fixed $L$ and $z$, $\beta\mapsto G_L(z;\beta)$ is the restriction to
$(0,\infty)$ of a function meromorphic on $\C$ and real-analytic on
$\R$. If it does not vanish identically, its zero set in $(0,\infty)$ is
closed and discrete.
\end{lemma}

\begin{proof}
$Z_L(\beta)$, $\int\Phi_0e^{N\beta\sum\ReTr U_p}$ and $\int\Phi_0\Phi_ze^{
N\beta\sum\ReTr U_p}$ are integrals over a compact space of functions
entire in $\beta$ with a uniform bound on compact sets of $\beta$, hence
entire; $Z_L(\beta)>0$ for real $\beta$. The zeros of a real-analytic
function on an interval that is not identically zero are isolated.
\end{proof}

The second ingredient is the strong-coupling expansion, which identifies
the leading Taylor coefficient. We recall the polymer representation in
the form of Osterwalder and Seiler~\cite{OSei78} and Seiler~\cite{Sei82}.
From now on $\rho$ is irreducible, so that $\chi$ and $\bar\chi$ are the
only irreducible characters occurring in $\chi+\bar\chi$; $\delta=1$ if
$\chi=\bar\chi$. For an irreducible $R\ne1$ let $n(R)$ be the least $n$
with $m^{(n)}_R>0$; then $n(\rho)=n(\bar\rho)=1$, $n(R)\ge2$ otherwise,
and $a_R(\beta)=c_R(\beta)/c_0(\beta)$ satisfies
\begin{equation}\label{rEtwo:eqactivity}
\begin{aligned}
a_R(\beta)&=O(\beta^{n(R)}),\qquad a_\rho(\beta)=\tfrac{1+\delta}2N\beta+
O(\beta^2),\\
c_0(\beta)&=1+O(\beta^2),
\end{aligned}
\end{equation}
the last because $\int\ReTr=0$ for a nontrivial irreducible $\rho$. For a
plaquette $p$ and $s\in\C$ let $c^{(p)}_R(\beta,s)$ be the character
coefficients of $e^{sF_p}e^{N\beta\ReTr U_p}=e^{sN}e^{(N\beta-s)\ReTr U_p}$
and $a^{(p)}_R(s)=c^{(p)}_R(\beta,s)/c^{(p)}_0(\beta,s)$; then $a^{(p)}_R(0)
=a_R(\beta)$ and, since $F_pe^{N\beta\ReTr}=\sum_n(N\beta)^n(N-\ReTr)
(\ReTr)^n/n!$,
\begin{equation}\label{rEtwo:eqsource}
\partial_sa^{(p)}_R(0)=O(\beta^{n(R)-1}),\qquad\partial_sa^{(p)}_\rho(0)=
-\tfrac{1+\delta}2+O(\beta).
\end{equation}
For complex $|s|\le r_N:=1/(2N)$ one has $|s\ReTr g|\le\frac12$, so
$\operatorname{Re}e^{-s\ReTr g}\ge e^{-1/2}\cos\frac12>\frac12$ and
$|c^{(p)}_0(\beta,s)|\ge\frac12e^{-1/2}c_0(\beta)$, while
$\sum_Rd_R|c^{(p)}_R(\beta,s)|\le e^{N/2}\sum_n((N\beta+\frac12)/2)^n(2N)^n/
n!=e^{N^2\beta+N}$, because $c^{(p)}_R(\beta,s)=e^{sN}\sum_n((N\beta-s)/2)^n
m^{(n)}_R/n!$ and $\sum_Rd_Rm^{(n)}_R=(2N)^n$; hence
\begin{equation}\label{rEtwo:eqsourcebound}
\sum_{R\ne1}d_R\,|a^{(p)}_R(s)|\le C_N\quad(|s|\le r_N,\ \beta\le1/N^2),
\end{equation}
and, by the same count restricted to $n\ge n(R)$, for $\beta\le1/N^2$,
\[
\sum_{R:\,n(R)\ge n}d_R\,a_R(\beta)\le2e\,(N^2\beta)^n.
\]

\begin{lemma}[polymer expansion]\label{rEtwo:polymer}
A \emph{polymer} is a pair $\gamma=(A,R)$ with $A$ a nonempty finite set
of plaquettes of $\Lambda_L$ that is connected under the relation of
sharing a link, and $R$ an assignment of a nontrivial irreducible
representation $R_a$ to each $a\in A$, such that $I(\gamma)=\int\prod_{a
\in A}\chi_{R_a}(U_a)\,dU\ne0$. Two polymers are compatible if their
plaquette sets share no link. For sources $s,s'$ at two plaquettes $p\ne
q$ let $W(\gamma;s,s')=I(\gamma)\prod_{a\in A}a^{(a)}_{R_a}$ with
$a^{(a)}_R=a^{(p)}_R(s)$, $a^{(q)}_R(s')$, $a_R(\beta)$ for $a=p$, $a=q$
and otherwise. Then:

\textup{(i)} If a link lies in exactly one plaquette of $A$, then
$I(A,R)=0$; in particular $|A|\ge6$ for every polymer. For every $(A,R)$,
$I(A,R)$ is the product of the $I$ of the link-connected components of $A$.

\textup{(ii)} $\int e^{sF_p+s'F_q}e^{N\beta\sum\ReTr U_a}dU=\prod_ac^{(a)}_0
\sum\prod_iW(\gamma_i;s,s')$, the sum over finite families of pairwise
compatible polymers, including the empty family.

\textup{(iii)} There is $\beta_0=\beta_0(N)>0$ such that for $0\le\beta
\le\beta_0$ and complex $|s|,|s'|\le r_N$ the Mayer expansion
\[
\log\frac{\int e^{sF_p+s'F_q}e^{N\beta\sum\ReTr U_a}dU}{\prod_ac^{(a)}_0}=
\sum_X\varphi(X)\prod_{\gamma\in X}W(\gamma;s,s')
\]
over clusters $X$ (finite multisets of polymers whose incompatibility
graph is connected, $\varphi$ the Ursell coefficients, $\varphi(X)=1$ for
one polymer) converges absolutely, uniformly in $L$, $s$, $s'$, and the
sum of $|\varphi(X)\prod W|$ over the clusters $X$ with $\sum_{\gamma\in X}
\sum_{a\in A_\gamma\setminus\{p,q\}}n(R_a)\ge n$ and with some polymer
containing $p$ is at most $K_N(C_N\beta)^n$, with constants depending on
$N$ only.
\end{lemma}

\begin{proof}
(i) If the link $e$ lies only in $a\in A$, write $U_a=XU_e^{\pm1}Y$ with
$X,Y$ independent of $U_e$; then
\begin{align*}
\chi_{R_a}(XU_e^{\pm1}Y)&=\sum_{i,j,k}R_a(X)_{ij}R_a(U_e^{\pm1})_{jk}R_a(Y)_{ki},\\
\int_G R_a(U)_{jk}\,dU&=0
\end{align*}
for nontrivial irreducible $R_a$. A set of plaquettes in which every link lies in at
least two plaquettes has at least six members, the faces of a cube being
the least case. The factorisation is Fubini. (ii) Expand each plaquette
factor as $c^{(a)}_0(1+\sum_{R\ne1}a^{(a)}_R\chi_R(U_a))$, multiply out,
integrate, and group the plaquettes carrying a nontrivial representation
into link-connected components; by (i) only the components with
$I\ne0$ survive, and distinct components share no link. (iii) is the
convergent strong-coupling cluster expansion of~\cite{OSei78} (see
also~\cite[Ch.~3]{Sei82}), with the sources carried along; we indicate the
bounds. Since $|\chi_R|\le d_R$, $|W(\gamma;s,s')|\le\prod_{a\in A}d_{R_a}
|a^{(a)}_{R_a}|$, and summing over the assignments $R$ with
\eqref{rEtwo:eqsourcebound} and the bound after it, the polymers with a
given plaquette set $A$ have total weight at most $C_N^{2}(2eN^2\beta)^{|A|
-2}$; the number of link-connected sets of $n$ plaquettes containing a
given link is at most $C^n$ with an absolute $C$, and $|A|\ge6$ by (i).
Hence for every polymer $\gamma_0$ the polymers incompatible with it,
which contain one of its at most $4|A_{\gamma_0}|$ links, satisfy
$\sum_{\gamma\not\sim\gamma_0}|W(\gamma)|e^{|A_\gamma|}\le|A_{\gamma_0}|$
once $\beta\le\beta_0(N)$, which is the Koteck\'y--Preiss
criterion~\cite{KP86} with $a(\gamma)=|A_\gamma|$, uniformly in $L$, $s$,
$s'$; it gives the absolute convergence of the Mayer series, uniformly on
the polydisc, and, applied with the activities of the plaquettes $a\ne p,q$
multiplied by $(\beta_0(N)/\beta)^{n(R_a)}$, whose sums are bounded by $2e\sum_k(N^2\beta_0(N))^k$ by the bound after \eqref{rEtwo:eqsourcebound}, the bound
$K_N(C_N\beta)^n$ on the clusters of order at least $n$ through $p$.
\end{proof}

\begin{lemma}[the minimal surfaces]\label{rEtwo:plateau}
Let $L\ge6$ even, $2\le t<L/2$, let $p$ be a plaquette with the corner $0$
and $q$ a plaquette with the corner $te_0$, and let $\Sigma$ be a finite
set of plaquettes of $\Lambda_L$ containing $p$ and $q$, connected under
the relation of sharing a link, such that every link of every plaquette
of $\Sigma$ lies in at least two plaquettes of $\Sigma$. Then $|\Sigma
\setminus\{p,q\}|\ge4t$, with equality if and only if $\Sigma=\partial
\mathcal B$ is the set of the $4t+2$ faces of a box $\mathcal B=[0,t]
\times I\times J$, where $I=[a,a+1]$ and $J=[b,b+1]$ lie in two distinct
spatial directions, $a,b\in\{-1,0\}$, and the fourth coordinate is $0$.
There are $12$ such boxes; each has exactly three faces with the corner
$0$ and three with the corner $te_0$, so there are $108$ triples
$(\mathcal B,p,q)$ with $p,q\in\partial\mathcal B$.
\end{lemma}

\begin{proof}
Identify the times with $\{-L/2+1,\dots,L/2\}$; the plaquettes with the
corner $0$ have their vertices at times in $\{-1,0,1\}$, those with the
corner $te_0$ at times in $\{t-1,t,t+1\}$. For a half-integer $\tau$ let
$E_\tau$ be the set of temporal plaquettes of $\Sigma$ with time extent
$[\tau-\tfrac12,\tau+\tfrac12]$; each is $[\tau-\tfrac12,\tau+\tfrac12]
\times e$ for a spatial link $e$, and we identify $E_\tau$ with a set of
spatial links. If a spatial site $y$ is an endpoint of a link of
$E_\tau$, the temporal link at $y$ over $[\tau-\tfrac12,\tau+\tfrac12]$
lies in a plaquette of $\Sigma$, hence in at least two, which are
plaquettes of $E_\tau$; so every vertex of the graph $E_\tau$ has degree
$0$ or at least $2$, and $E_\tau$, if nonempty, contains a cycle of the
graph $\Lambda^3$, which has at least $4$ links since $L\ge6$: cycles of
$\Lambda^3$ have even length, and a cycle of length $4$ is the boundary
of a unit square because a straight closed path needs length $L$.

Every plaquette has its vertices at one time or at two consecutive times,
and two plaquettes sharing a link have a common time. Along a chain of
plaquettes of $\Sigma$ from $p$ to $q$, consecutive members sharing a
link, the union of the times of the members is therefore an arc of the
time circle containing $0$ and $t$, and the passage from a time $s$ to
$s+1$ inside the arc is made by a member with vertices at both times, that
is, by a temporal plaquette over $[s,s+1]$. Let $s\in\{0,\dots,t-1\}$. If
$E_{s+1/2}=\emptyset$, the arc does not contain the step from $s$ to
$s+1$, so it runs from $0$ through $-1,-2,\dots$ to $t$, and $E_\tau\ne
\emptyset$ for each of the $L-t$ half-integers $\tau$ between $t$ and
$L\equiv0$; then $|\Sigma\setminus\{p,q\}|\ge4(L-t)-2>4t$ because $L-t\ge
t+1$. So assume $E_{s+1/2}\ne\emptyset$, hence $|E_{s+1/2}|\ge4$, for all
$s\in\{0,\dots,t-1\}$. The plaquettes $p$ and $q$ belong to at most one
of these sets each, $p$ to $E_{1/2}$ if $p$ is temporal over $[0,1]$, $q$
to $E_{t-1/2}$ if $q$ is temporal over $[t-1,t]$. If $p$ is temporal over
$[0,1]$, its spatial link at time $0$ lies in a second plaquette of
$\Sigma$, which is a spatial plaquette at time $0$ or a temporal plaquette
over $[-1,0]$ (another temporal plaquette over $[0,1]$ through the same
spatial link would be $p$ itself); in the second case $|E_{-1/2}|\ge4$. In
every case $\Sigma\setminus\{p,q\}$ contains the $\sum_s|E_{s+1/2}|-[p\in
E_{1/2}]-[q\in E_{t-1/2}]\ge4t-2$ plaquettes of the $E_{s+1/2}$, $0\le s
\le t-1$, other than $p,q$, plus at least one further plaquette, at times
in $\{-1,0\}$, for $p$ if $p$ is temporal, and one at times in $\{t,t+1\}$
for $q$ if $q$ is temporal; these are distinct since $t\ge2$. If $p$ is
spatial, it is not counted in the $E$'s and the bound $4t$ holds
directly. The bound $4t$ follows.

Equality forces: $|E_{s+1/2}|=4$ for $0\le s\le t-1$, so $E_{s+1/2}=
\partial Q_s$ for a unit square $Q_s$ of $\Lambda^3$; $E_\tau=\emptyset$
for every other $\tau$ (a nonempty $E_\tau$ has at least four members, at
most one of which is $p$ or $q$, and would add at least three plaquettes
to the count); and no spatial plaquettes besides $p$, $q$ and one spatial
plaquette at time $0$ (resp.\ $t$) when $p$ (resp.\ $q$) is temporal. For
$0\le s\le t-2$ the four spatial links $\partial Q_s$ at time $s+1$ lie in
the plaquettes of $E_{s+1/2}$ and need a second plaquette of $\Sigma$
each; spatial plaquettes at time $s+1\in[1,t-1]$ are excluded, so these
links lie in $E_{s+3/2}=\partial Q_{s+1}$, and $Q_{s+1}=Q_s=:Q$. The links
$\partial Q$ at time $0$ need a second plaquette, which is spatial at time
$0$ since $E_{-1/2}=\emptyset$: if $p$ is spatial it is $p$, so $p=Q\times
\{0\}$ and $0\in Q$; if $p$ is temporal, $p\in E_{1/2}$ and the single
allowed spatial plaquette at time $0$ must contain all four links, so it
is $Q\times\{0\}$, and $0\in p$ gives $0\in Q$ again. Symmetrically at
time $t$. Hence $\Sigma$ consists of the tube $\bigcup_s\partial Q\times
[s,s+1]$ and the two faces $Q\times\{0\}$, $Q\times\{t\}$, which is
$\partial\mathcal B$ for $\mathcal B=[0,t]\times Q$ with $Q=I\times J\ni0$
as stated; conversely each such $\partial\mathcal B$ satisfies the
hypotheses with equality. The unit squares $Q\ni0$ in the three spatial
coordinate planes number $3\cdot4=12$. The faces of $\partial\mathcal B$
containing the vertex $(0,0)$ are the end face $Q\times\{0\}$ and the two
side faces $e\times[0,1]$ for the two links $e$ of $\partial Q$ at $0$;
likewise at $(t,0)$.
\end{proof}

\begin{lemma}[weight of a box]\label{rEtwo:boxweight}
Let $\mathcal B$ be as in Lemma~\textup{\ref{rEtwo:plateau}} and let
$R$ assign $\rho$ or $\bar\rho$ to each face of $\partial\mathcal B$. Then
$I(\partial\mathcal B,R)=N^{-4t}$ if the assignment is \emph{consistent},
that is, if for every link the two faces containing it carry mutually
dual representations with respect to the orientations induced by the
plaquette variables, and $I=0$ otherwise. There are exactly $2-\delta$
consistent assignments. If some face carries a representation $R_a\notin
\{\rho,\bar\rho\}$ while its neighbours carry $\rho$ or $\bar\rho$, then
$I=0$.
\end{lemma}

\begin{proof}
Each link $e$ of $\partial\mathcal B$ lies in exactly two faces. Writing
the two plaquette variables as $XU_e^{\epsilon}Y$ and $X'U_e^{\epsilon'}Y'$,
the integral over $U_e$ of the product of the two matrix coefficients of
$U_e$ vanishes unless the two representations of $U_e$ appearing are dual
to each other, that is, unless $R'$ restricted along $e$ is $\bar R$ with
respect to the induced orientations; when it does not vanish it equals
$\delta\delta/N$ by the Schur relations. This gives the last statement.
Choosing the orientation of $\partial\mathcal B$ as the boundary of the
box, each link is traversed in opposite directions by its two faces, so a
consistent assignment is one in which every face carries $\rho$ with
respect to the boundary orientation, or every face carries $\bar\rho$;
these two coincide when $\chi=\bar\chi$. For a consistent assignment the
integral is the partition function of two-dimensional lattice gauge theory
on the sphere $\partial\mathcal B$ with the face weight $\chi_\rho$: after
the link integrals, the remaining index contractions close up into one
loop at each of the $4(t+1)$ vertices, giving $N^{V}$, against the factor
$N^{-1}$ of each of the $8t+4$ links, and $V-E=-4t$.
\end{proof}

\begin{theorem}[leading strong-coupling term]\label{rEtwo:taylor}
Let $\rho$ be irreducible, $L\ge6$ even and $2\le t<L/2$. There are
$\beta_0(N)>0$ and $C(N,t)<\infty$, independent of $L$, such that for
$0\le\beta\le\beta_0(N)$
\[
\Bigl|G_L(te_0;\beta)-108\Bigl(\frac{1+\delta}2\Bigr)^{4t+1}\beta^{4t}\Bigr|
\le C(N,t)\,\beta^{4t+1}.
\]
The same holds for $G(te_0)$ in every torus limit.
\end{theorem}

\begin{proof}
Since the left side is the second mixed cumulant,
\[
G_L(te_0)=\sum_{p\ni0,\,q\ni te_0}\partial_s\partial_{s'}\log\int e^{sF_p
+s'F_q}e^{N\beta\sum\ReTr}\Big|_{s=s'=0},
\]
a sum of $576$ terms. Apply Lemma~\ref{rEtwo:polymer}(iii) to each term:
the series converges uniformly on the polydisc $|s|,|s'|\le r_N$, so it
may be differentiated term by term. The volume factors $\prod_ac^{(a)}_0$
contribute $\log c^{(p)}_0(s)+\log c^{(q)}_0(s')$, whose mixed derivative
is $0$; a cluster contributes a nonzero mixed derivative only if some
polymer of it contains $p$ and some polymer contains $q$. For such a
cluster $X$ let $\Sigma$ be the union of the plaquette sets of its
polymers. It contains $p$ and $q$, is connected under sharing a link (each
polymer is, and incompatible polymers share a link), and every link of a
plaquette of $\Sigma$ lies in at least two plaquettes of the polymer
containing that plaquette, by Lemma~\ref{rEtwo:polymer}(i). By
Lemma~\ref{rEtwo:plateau}, $\sum_{\gamma\in X}\sum_{a\in A_\gamma\setminus
\{p,q\}}n(R_a)\ge|\Sigma\setminus\{p,q\}|\ge4t$, with equality only if
$\Sigma=\partial\mathcal B$, every $n(R_a)=1$ for $a\ne p,q$, and no
plaquette other than $p,q$ belongs to two polymers of $X$. In the equality
case $X$ consists of a single polymer: a polymer whose plaquettes lie in
$\{p,q\}$ does not exist (a polymer has at
least six plaquettes), and if two polymers of $X$ partitioned the faces of
the sphere $\partial\mathcal B$, a link whose two faces lie in different
polymers would lie in one plaquette of each, contradicting
Lemma~\ref{rEtwo:polymer}(i); moreover $R_p,R_q\in\{\rho,\bar\rho\}$ by the
last statement of Lemma~\ref{rEtwo:boxweight}. By
Lemma~\ref{rEtwo:polymer}(iii) and the Cauchy estimates on the polydisc,
the clusters of order at least $4t+1$ contribute at most $r_N^{-2}K_N(C_N
\beta)^{4t+1}$ to the mixed derivative of each of the $576$ terms. The
order-$4t$ part is the sum over the $108$ triples $(\mathcal B,p,q)$ of
$\partial_s\partial_{s'}W(\partial\mathcal B,R;s,s')$ summed over the
consistent $R$ of Lemma~\ref{rEtwo:boxweight}, which by
\eqref{rEtwo:eqactivity}, \eqref{rEtwo:eqsource} equals
\begin{align*}
&(2-\delta)\,N^{-4t}\Bigl(\frac{1+\delta}2N\beta\Bigr)^{4t}\Bigl(\frac{1+
\delta}2\Bigr)^2\bigl(1+O(\beta)\bigr)\\
&\qquad=\Bigl(\frac{1+\delta}2\Bigr)^{4t+1}\beta^{4t}\bigl(1+O(\beta)\bigr),
\end{align*}
using $(2-\delta)(1+\delta)/2=1$ for $\delta\in\{0,1\}$. The $O(\beta)$ is
uniform in $L$. The statement for torus limits follows since $G(te_0)=
\lim G_L(te_0)$ along the subsequence.
\end{proof}

\begin{corollary}[the exceptional set]\label{rEtwo:exceptional}
Let $\rho$ be irreducible and $L\ge6$ even. The set $\mathcal D_L=\{\beta>0:
G_L(2e_0;\beta)=0\}$ is closed and discrete in $(0,\infty)$, and there is
$\beta_1(N)>0$ independent of $L$ with $\mathcal D_L\cap(0,\beta_1(N)]=
\emptyset$. For $\beta\notin\mathcal D_L$, $G_L(te_\mu;\beta)>0$ for all
$\mu$ and all $2\le t\le L-2$; for $\beta\in\mathcal D_L$ the degenerate
alternative of Theorem~\textup{\ref{rEtwo:axial}(iii),(iv)} holds. For
every torus limit at $\beta\in(0,\beta_1(N)]$, $G(te_0)>0$ for all
$t\ge2$.
\end{corollary}

\begin{proof}
Put
\[
\beta_1(N)=\min\Bigl(\beta_0(N),\ \min_{t=2,3,4}54\bigl(\tfrac{1+\delta}2
\bigr)^{4t+1}/C(N,t)\Bigr).
\]
By Theorem~\ref{rEtwo:taylor} with $t=2$,
$G_L(2e_0;\beta)\ge\frac12\cdot108(\frac{1+\delta}2)^9\beta^8>0$ for $0<\beta
\le\beta_1(N)$, so $\beta\mapsto G_L(2e_0;\beta)$ does not vanish
identically and Lemma~\ref{rEtwo:analytic} applies; the rest of the
finite-torus statement is Theorem~\ref{rEtwo:axial}(iii),(iv). For a torus
limit at $\beta\le\beta_1(N)$, Theorem~\ref{rEtwo:taylor} with $t=2,3,4$
(admissible since $L\to\infty$ along the subsequence) gives $G(te_0)>0$
for $t=2,3,4$, so the first alternative of Corollary~\ref{rEtwo:limit}(ii)
holds and $G(te_0)>0$ for all $t\ge2$.
\end{proof}

\begin{proposition}[sum rule]\label{rEtwo:sumrule}
With $S=\sum_{p}F_p$ the Wilson action of $\Lambda_L$, for every $\beta>0$
\[
\sum_{z\in\Lambda_L}G_L(z)=\frac{16}{L^4}\Var_L(S)=-\frac{16}
{NL^4}\,\frac{d}{d\beta}\langle S\rangle_L>0 .
\]
\end{proposition}

\begin{proof}
Every plaquette has four corners, so $\sum_x\Phi_x=4S$ and
\[
\sum_zG_L(z)=L^{-4}\sum_{x,z}\Cov(\Phi_x,\Phi_{x+z})=L^{-4}\Var(4S).
\]
The density is
$e^{-N\beta S}$ up to a constant, so $d\langle S\rangle_L/d\beta=-N\Var_L(S)$,
and $\Var_L(S)>0$ because $S$ is a nonconstant continuous function and the
Wilson measure has a strictly positive density.
\end{proof}

\begin{remark}[what is and is not proved about nonvanishing]
\label{rEtwo:status1}
For $\Phi$ on a finite torus, the connected correlation at an axial
separation $te_\mu$, $2\le t\le L-2$, is strictly positive for every
coupling outside the closed discrete set $\mathcal D_L$, which is empty
in $(0,\beta_1(N)]$, and its vanishing at one coupling is equivalent to
the operator identity $B=c_LT^2$ on $\mathfrak h_P$, which forces the
vanishing at every axial separation and at every $z$ with a coordinate of
modulus in $[2,L-2]$, and forces the one-point function to be independent
of the time extent. Whether $\mathcal D_L$ is empty for every $L$ is not
decided here; it is decided for a single spatial plaquette
(Theorem~\ref{rEtwo:plaquette}). For non-axial $z$ only the
real-analyticity of $\beta\mapsto G_L(z;\beta)$ (Lemma~\ref{rEtwo:analytic})
and the vanishing in the degenerate alternative are obtained; neither the
sign of $G_L(z;\beta)$ nor the sign of its first nonvanishing Taylor
coefficient is asserted. The sum rule shows that $G_L$ does not vanish
identically in $z$ at any coupling.
\end{remark}

\subsubsection{Quantitative lower bounds along an axis}

The dichotomy of Theorem~\ref{rEtwo:axial} comes with a rate. In the
nondegenerate alternative put, for $\beta\notin\mathcal D_L$,
\[
m_L(t;\beta)=\log\frac{G_L(te_0;\beta)}{G_L((t+1)e_0;\beta)}\qquad
(2\le t\le L-3),
\]
and for a torus limit in the first alternative of
Corollary~\ref{rEtwo:limit}(ii) let $m(t)$ be as in
Corollary~\ref{rEtwo:limit}(iii).

\begin{theorem}[secant bound]\label{rEtwo:logconvex}
Let $\beta>0$, $L\ge6$ even and $\beta\notin\mathcal D_L$. Then $t\mapsto
m_L(t;\beta)$ is nonincreasing on $\{2,\dots,L-3\}$, $m_L(t;\beta)\ge0$ for
$t\le L/2-1$, $m_L(L-1-t;\beta)=-m_L(t;\beta)$, and for $2\le t_0<t\le L-2$
\begin{align*}
G_L(te_0;\beta)&\ge G_L(t_0e_0;\beta)\,e^{-(t-t_0)\,m_L(t_0;\beta)}\\
&\ge G_L(t_0e_0;\beta)\,e^{-(t-t_0)\,m_L(2;\beta)} .
\end{align*}
For a torus limit $\omega$ at $\beta>0$ in the first alternative of
Corollary~\textup{\ref{rEtwo:limit}(ii)}, and all $2\le t_0<t$,
\[
G(te_0)\ge G(t_0e_0)\,e^{-(t-t_0)m(t_0)}\ge G(t_0e_0)\,e^{-(t-t_0)m(2)},
\]
and $m(t_0)\downarrow m_\Phi$, so that the exponential rate of the bound
improves to $m_\Phi$ as $t_0\to\infty$; moreover
\[
G(te_0)\le G(t_0e_0)\,e^{-(t-t_0)m(t-1)},\quad
\lim_{t\to\infty}\frac{\log G(te_0)}{t}=-m_\Phi,
\]
so $m_\Phi$ is the exact exponential rate of $G$ along the axis.
\end{theorem}

\begin{proof}
By Corollary~\ref{rEtwo:exceptional}, $G_L(te_0;\beta)>0$ for $2\le t\le
L-2$, so $g(t)=\log G_L(te_0;\beta)$ is defined there, and
Theorem~\ref{rEtwo:axial}(ii) says that $g$ is convex on $\{2,\dots,L-2\}$:
its increments $g(t+1)-g(t)=-m_L(t;\beta)$ are nondecreasing in $t$. The
symmetry $g(t)=g(L-t)$ of Theorem~\ref{rEtwo:axial}(i) gives $m_L(L-1-t)=
-m_L(t)$, and the monotonicity of Theorem~\ref{rEtwo:axial}(ii) gives
$m_L(t)\ge0$ for $t+1\le L/2$. For $t_0<t$,
\begin{align*}
g(t)-g(t_0)&=\sum_{s=t_0}^{t-1}\bigl(g(s+1)-g(s)\bigr)\\
&\ge(t-t_0)\bigl(g(t_0+1)-g(t_0)\bigr),
\end{align*}
which is the first inequality; the second follows from $m_L(t_0)\le
m_L(2)$. In a torus limit the same argument applies to $g(t)=\log G(te_0)$,
$t\ge2$, which is finite by the first alternative and the positivity
$G(2e_0)\ge G(3e_0)>0$ of Corollary~\ref{rEtwo:limit}(i),(ii); the
monotonicity and the limit $m_\Phi$ of the $m(t)$ are
Corollary~\ref{rEtwo:limit}(iii). The upper bound is the convexity
inequality $g(t)-g(t_0)\le(t-t_0)(g(t)-g(t-1))$. Finally $t^{-1}(g(t)-g(2))
=-t^{-1}\sum_{s=2}^{t-1}m(s)\to-m_\Phi$ by Ces\`aro.
\end{proof}

At strong coupling the two numbers entering the bound are explicit.

\begin{theorem}[explicit constants at strong coupling]\label{rEtwo:strong}
Let $\rho$ be irreducible. There are $\beta_2(N)\in(0,\beta_1(N)]$ and
$C_2(N)<\infty$ such that for $0<\beta\le\beta_2(N)$ and every even
$L\ge8$ the numbers
\[
c(\beta)=108\Bigl(\frac{1+\delta}2\Bigr)^{9}\beta^8,\qquad
m(\beta)=4\log\frac2{(1+\delta)\beta}
\]
satisfy
\[
\Bigl|\frac{G_L(2e_0;\beta)}{c(\beta)}-1\Bigr|\le C_2(N)\beta,\quad
|m_L(2;\beta)-m(\beta)|\le C_2(N)\beta ,
\]
and consequently, for $2\le t\le L-2$,
\begin{align*}
G_L(te_0;\beta)&\ge\bigl(1-C_2(N)\beta\bigr)\,c(\beta)\,
e^{-(t-2)\,(m(\beta)+C_2(N)\beta)}\\
&=\bigl(1-C_2(N)\beta\bigr)\,108\Bigl(\frac{1+\delta}2\Bigr)^{9}\beta^{8}\\
&\qquad\times\Bigl(\frac{(1+\delta)\beta}{2}\Bigr)^{4(t-2)}e^{-C_2(N)\beta(t-2)} .
\end{align*}
The same inequalities hold for $G(te_0)$ and $m(2)$ in every torus limit
at $\beta\in(0,\beta_2(N)]$, for all $t\ge2$. The constants $\beta_2(N)$
and $C_2(N)$ do not depend on $L$.
\end{theorem}

\begin{proof}
Theorem~\ref{rEtwo:taylor} with $t=2$ and $t=3$ (admissible since $3<L/2$
for $L\ge8$) gives, for $\beta\le\beta_0(N)$,
\begin{align*}
G_L(2e_0;\beta)&=108\Bigl(\frac{1+\delta}2\Bigr)^{9}\beta^8\bigl(1+
\varepsilon_2(\beta)\bigr),\\
G_L(3e_0;\beta)&=108\Bigl(\frac{1+\delta}2\Bigr)^{13}\beta^{12}\bigl(1+
\varepsilon_3(\beta)\bigr),
\end{align*}
with $|\varepsilon_t(\beta)|\le C(N,t)\beta/(108((1+\delta)/2)^{4t+1})$.
Choose $\beta_2(N)\le\beta_1(N)$ so small that $|\varepsilon_2|,
|\varepsilon_3|\le1/2$ on $(0,\beta_2(N)]$. Then both values are positive,
and
\[
m_L(2;\beta)=\log\frac{G_L(2e_0;\beta)}{G_L(3e_0;\beta)}=4\log\frac2{(1+
\delta)\beta}+\log\frac{1+\varepsilon_2(\beta)}{1+\varepsilon_3(\beta)},
\]
where
\[
\bigl|\log\bigl((1+\varepsilon_2)/(1+\varepsilon_3)\bigr)\bigr|
\le2(|\varepsilon_2|+|\varepsilon_3|)
\]
because $|\log(1+x)|\le2|x|$ for $|x|\le1/2$. This gives
the two estimates with a $C_2(N)$ depending only on $C(N,2)$, $C(N,3)$ and
$\delta$. The lower bound on $G_L(te_0;\beta)$ is then
Theorem~\ref{rEtwo:logconvex} with $t_0=2$, and $e^{-(t-2)m(\beta)}=
((1+\delta)\beta/2)^{4(t-2)}$. For torus limits, Theorem~\ref{rEtwo:taylor}
holds for $G(te_0)$ as stated, the first alternative holds by
Corollary~\ref{rEtwo:exceptional}, and the bound is
Theorem~\ref{rEtwo:logconvex}.
\end{proof}

\begin{remark}\label{rEtwo:strongremark}
(a) The bound of Theorem~\ref{rEtwo:strong} is sharp to leading order at
$t=2,3$ and, by Theorem~\ref{rEtwo:taylor}, sharp in the exponent at every
fixed $t<L/2$: $G_L(te_0;\beta)=108((1+\delta)/2)^{4t+1}\beta^{4t}(1+O(
\beta))$, while the bound gives $108((1+\delta)/2)^{4t+1}\beta^{4t}(1-O(
t\beta))$. The factor $e^{-C_2(N)\beta(t-2)}$ is the price of using only
the two shortest distances; it is removed, at any fixed $t_0$, by starting
the secant at $t_0$ instead of $2$, since $m_L(t_0;\beta)=4\log(2/((1+
\delta)\beta))+O(\beta)$ for every fixed $t_0<L/2-1$ by the same
computation.

(b) In the language of Sec.~\ref{ssec:limitE}, the rate is a lattice mass.
Since $m(t)\downarrow m_\Phi$, Theorem~\ref{rEtwo:strong} gives $m_\Phi\le
m(2)=4\log(2/((1+\delta)\beta))+O(\beta)$ in every torus limit at $\beta
\le\beta_2(N)$, and by (a) the same upper bound holds with $m(t_0)$ for
every fixed $t_0$. The equality $m_\Phi=4\log(2/((1+\delta)\beta))+O(\beta)$
would require the remainder $C(N,t)\beta^{4t+1}$ of
Theorem~\ref{rEtwo:taylor} to be of the form $C(N)^t\beta^{4t+1}$ uniformly
in $t$; this is not established here. The number $4\log(1/\beta)$ is the
area-law cost of the four side plaquettes per unit length of the tube of
Lemma~\ref{rEtwo:plateau}; in the normalisation $\beta_W=N^2\beta$ of
Sec.~\ref{sec:intro} the bound reads $m_\Phi\le4\log(2N^2/((1+\delta)
\beta_W))+O(\beta_W/N^2)$, the leading strong-coupling estimate of the
lightest glueball mass in the plaquette channel, for which see~\cite{MM94}.
\end{remark}

\paragraph{The natural parameter and the degeneration of the bound.}
Let $\nu_\beta$ be the one-link measure $c_0(\beta)^{-1}e^{N\beta\ReTr g}
\,dg$ on $G$ and
\[
u(\beta)=\frac{a_\rho(\beta)}N=\frac1N\int_G\chi(g)\,d\nu_\beta(g)
=\frac1N\int_G\ReTr g\,d\nu_\beta(g),
\]
the second equality because $g\mapsto g^{-1}$ preserves $\nu_\beta$ and
$\chi(g^{-1})=\overline{\chi(g)}$. Then $u$ is real-analytic on $[0,
\infty)$, $u(\beta)=\tfrac{1+\delta}2\beta+O(\beta^2)$ by
\eqref{rEtwo:eqactivity}, $u'(\beta)=\Var_{\nu_\beta}(\ReTr g)>0$ since
$\ReTr g$ is not constant on $G\ne\{1\}$, and $u(\beta)\uparrow1$ as
$\beta\to\infty$: writing $\ReTr g=N-\tfrac12\|g-1\|_{\mathrm{HS}}^2$, the
set $\{\ReTr g<N-\varepsilon\}$ has $\nu_\beta$-measure at most
$e^{-N\beta\varepsilon/2}/\operatorname{Haar}\{\ReTr g\ge N-\varepsilon/2\}
\to0$. The leading strong-coupling term of Theorem~\ref{rEtwo:taylor} is
$108\,\tfrac{1+\delta}2\,u(\beta)^{4t}(1+O(\beta))$, each of the $4t$
plaquettes of the box contributing one factor $u$ (the activity $a_\rho$
divided by the dimension $N$ of Lemma~\ref{rEtwo:boxweight}), so that
$m(\beta)=4\log(1/u(\beta))+O(\beta)$. As $\beta$ grows, $u(\beta)$
increases to $1$ and the formal rate $4\log(1/u(\beta))$ decreases to $0$.
What is proved is the following.

\begin{proposition}[the bound at every coupling]\label{rEtwo:degenerate}
Let $\rho$ be irreducible and $L\ge8$ even.

\textup{(i)} For every $\beta\in(0,\infty)\setminus\mathcal D_L$ and
$2\le t\le L-2$,
\[
G_L(te_0;\beta)\ge G_L(2e_0;\beta)\,e^{-(t-2)m_L(2;\beta)},
\]
where $\beta\mapsto G_L(2e_0;\beta)>0$ and $\beta\mapsto
m_L(2;\beta)\ge0$ are real-analytic on $(0,\infty)\setminus\mathcal D_L$.
For $\beta\in\mathcal D_L$ both sides vanish.

\textup{(ii)} On $(0,\beta_2(N)]$ the two functions are given by
Theorem~\ref{rEtwo:strong}; in particular $m_L(2;\beta)\to\infty$ as
$\beta\downarrow0$: the exponential rate $4\log(2/((1+\delta)\beta))$ of
the bound diverges while its prefactor $G_L(2e_0;\beta)\asymp\beta^8$
vanishes. For $\beta>\beta_2(N)$ no bound on $G_L(2e_0;\beta)$ from below
or on $m_L(2;\beta)$ from above that is uniform in $L$ is obtained here:
Lemma~\ref{rEtwo:polymer}(iii) controls the Taylor series only for $\beta
\le\beta_0(N)$, and Lemma~\ref{rEtwo:analytic} gives analyticity in
$\beta$ for each $L$ without bounds uniform in $L$.

\textup{(iii)} If a sequence $\omega_k$ of torus limits at couplings
$\beta_k$ satisfies (AF) with constants $c_1\le C_1$, $\gamma$, $r_0$, then
for all $k\ge k_0$ with $4a_k\le r_0$
\[
\Bigl|m_k(3)-8\log\frac43\Bigr|\le\log\frac{C_1}{c_1}+\gamma\log
\frac{\ell(3a_k)}{\ell(4a_k)} ,
\]
where $m_k(3)=\log(G_k(3e_0)/G_k(4e_0))$ and the last term tends to $0$
as $a_k\to0$. In particular the lattice rates $m_k(3)$ stay in a bounded
window, whereas the strong-coupling rate $4\log(2/((1+\delta)\beta_k))$
leaves every bounded window as $\beta_k\to0$ and the formal rate
$4\log(1/u(\beta_k))$ tends to $0$ as $\beta_k\to\infty$.
\end{proposition}

\begin{proof}
(i) is Theorem~\ref{rEtwo:logconvex} with $t_0=2$ together with
Lemma~\ref{rEtwo:analytic} (the ratio of two positive real-analytic
functions is real-analytic, and $m_L(2;\beta)\ge0$ since $2\le L/2-1$) and
Theorem~\ref{rEtwo:axial}(iii). (ii) is Theorem~\ref{rEtwo:strong} and a
statement about the proofs given. For (iii), (AF) applied at $z=3e_0$ and
$z=4e_0$, both with $|z|_\infty\ge3$ and $a_k|z|\le r_0$, gives
\[
\frac{c_1w_\gamma(3a_k)}{C_1w_\gamma(4a_k)}\le\frac{G_k(3e_0)}{G_k(4e_0)}
\le\frac{C_1w_\gamma(3a_k)}{c_1w_\gamma(4a_k)},
\]
and $w_\gamma(3a_k)/w_\gamma(4a_k)=(4/3)^8(\ell(4a_k)/\ell(3a_k))^\gamma$,
with $1\le\ell(3a)/\ell(4a)\to1$ as $a\to0$.
\end{proof}

\paragraph{What reflection positivity gives.} The question whether
reflection positivity alone yields a lower bound along a reflection axis
has a precise answer.

\begin{proposition}[reflection positivity and lower bounds]
\label{rEtwo:rp}
Let $\omega$ be a torus limit at $\beta\ge0$, $\vartheta$ the reflection
through the hyperplane $x_0=0$ and $\vartheta'$ the reflection through
$x_0=\tfrac12$ of Lemma~\textup{\ref{rE:torus}(iii)}.

\textup{(i)} Reflection positivity gives $\omega((\Phi_{se_0}-c)\,
\vartheta(\Phi_{se_0}-c))=G(2se_0)\ge0$ and $\omega((\Phi_{(s+1)e_0}-c)\,
\vartheta'(\Phi_{(s+1)e_0}-c))=G((2s+1)e_0)\ge0$ for $s\ge1$, hence
$G(te_0)\ge0$ for all $t\ge2$; with the transfer-matrix reconstruction it
gives the log-convexity and monotonicity of Corollary~\ref{rEtwo:limit}
and the secant bound of Theorem~\ref{rEtwo:logconvex}: a lower bound for
$G(te_0)$ at every $t$ in terms of the two values $G(t_0e_0)$, $G((t_0+1)
e_0)$.

\textup{(ii)} No positive lower bound on $G(te_0)$, $t\ge2$, follows from
reflection positivity, invariance and monotonicity alone: the function
$G\equiv0$ on $\{|z|_\infty\ge2\}$ satisfies every inequality that these
properties impose, and it is realised, with $c=\omega(\Phi_0)=24N$, by the
product Haar measure $\beta=0$, which is a reflection-positive, translation-
and hypercubically invariant torus limit. Within the Wilson measures at
$\beta>0$ the same conclusion holds at every $\beta\in\mathcal D_L$ on the
torus $\Lambda_L$, by Theorem~\textup{\ref{rEtwo:axial}(iv)}.

\textup{(iii)} Monotonicity gives $G(te_0)\ge\lim_{s\to\infty}G(se_0)=
\sigma(\{1\})$ for all $t\ge3$, the mass of the spectral measure of
Corollary~\ref{rEtwo:limit} at $\lambda=1$; this is $0$ whenever
$\gamma'<1$, in particular whenever $\omega$ has a mass gap for the
reflection $\vartheta'$, so monotonicity yields no positive lower bound in
the situations of interest.
\end{proposition}

\begin{proof}
(i) For $F=\Phi_{se_0}-c$, which depends on the links in $\{x_0\ge s-1\}
\subset\{x_0\ge0\}$ when $s\ge1$, reflection positivity for $\vartheta$
gives $\omega(\bar F\vartheta F)\ge0$, and $\vartheta\Phi_{se_0}=
\Phi_{-se_0}$ because $\vartheta$ maps the plaquettes with a corner at
$se_0$ onto those with a corner at $-se_0$ and $F_p$ is invariant under
orientation reversal; so $\omega(\bar F\vartheta F)=\omega(\Phi_{se_0}
\Phi_{-se_0})-c^2=G(2se_0)$ by invariance. For $\vartheta'$, $\Phi_{(s+1)
e_0}$ depends on links in $\{x_0\ge s\}\subset\{x_0\ge1\}$ for $s\ge1$,
and $\vartheta'\Phi_{(s+1)e_0}=\Phi_{-se_0}$, giving $G((2s+1)e_0)\ge0$.
The rest of (i) is Corollary~\ref{rEtwo:limit} and
Theorem~\ref{rEtwo:logconvex}, whose proofs use only the Osterwalder--Seiler
reconstruction from reflection positivity.

(ii) At $\beta=0$ the links are independent Haar variables. For
$|z|_\infty\ge2$ let $p\ni0$ and $q\ni z$ be plaquettes; then $p\ne q$,
since the corners of $p$ lie in $\{|x|_\infty\le1\}$, and two distinct
plaquettes share at most one link. If they share no link, $F_p$ and $F_q$
are independent. If they share the link $l$, then conditionally on $U_l$
the remaining three links of $p$ are independent Haar variables, and, with
the links of $p$ taken in cyclic order from $l$ and $P_0=\int_G\rho(g)\,dg$
the orthogonal projection onto the $G$-invariant vectors of $\C^N$,
\begin{align*}
E(F_p\mid U_l)&=N-\ReTr\bigl(U_l^{\pm1}E(U_{l_1}^{\pm1}U_{l_2}^{\pm1}U_{l_3}
^{\pm1})\bigr)\\
&=N-\ReTr(U_l^{\pm1}P_0^3)=N-\Tr P_0 ,
\end{align*}
because Haar measure is inversion invariant, $P_0^3=P_0$ and $\rho(g)P_0=
P_0$; this is a constant, equal to $E(F_p)$. Hence $E(F_pF_q)=E(E(F_p\mid
U_l)F_q)=E(F_p)E(F_q)$, since $F_q$ depends on $U_l$ and on links
independent of those of $p$, and summing over $p\ni0$, $q\ni z$ gives
$G(z)=0$. The product measure is reflection positive and invariant, and
$c=24(N-\Tr P_0)=24N$ for irreducible $\rho$. The statement about
$\mathcal D_L$ is Theorem~\ref{rEtwo:axial}(iv).

(iii) By Corollary~\ref{rEtwo:limit}(i), $G((3+n)e_0)=\int\lambda^n\,
d\sigma\downarrow\sigma(\{1\})$ by dominated convergence, and $\sigma$ is
carried by $[0,\gamma']$ because $\Psi'\perp\Omega'$, as shown in the proof
of Corollary~\ref{rEtwo:limit}(iii).
\end{proof}

\begin{remark}\label{rEtwo:rpremark}
Proposition~\ref{rEtwo:rp} separates the two inputs of a lower bound:
reflection positivity propagates positivity from two short distances to
all distances along the axis, with the exact rate $m_\Phi$ in the limit,
while the nonvanishing at short distance is a property of the specific
measure, established here by analyticity and the strong-coupling
expansion outside the discrete set $\mathcal D_L$
(Corollary~\ref{rEtwo:exceptional}) and with explicit values for
$\beta\le\beta_2(N)$ (Theorem~\ref{rEtwo:strong}). For a non-axial $z$
there is no reflection in a coordinate hyperplane mapping $0$ to $z$; the
diagonal reflections are not among the reflection positivities proved in
Lemma~\ref{rE:torus}(iii), and the positivity $G(z)\ge0$ for general $z$
is not asserted here.
\end{remark}

\subsubsection{Comparison with the perturbative law and the missing step}

Let $\omega_k$ be torus limits at couplings $\beta_k$ with spacings $a_k
\to0$, and $G_k$, $Z_k$, $c_1\le C_1$, $\gamma$, $r_0$, $\ell$, $w_\gamma$
as in (AF) of Sec.~\ref{ssec:limitE}; with $r=1$ the range of (AF) is
$|z|_\infty\ge3$, and for $z=te_0$ one has $|z|=t$. For $t\ge2$ put
$m_k(t)=\log(G_k(te_0)/G_k((t+1)e_0))$ when $G_k((t+1)e_0)>0$.
Theorem~\ref{rEtwo:logconvex} gives, in the first alternative of
Corollary~\ref{rEtwo:limit}(ii), the exponential lower bound
\begin{equation}\label{rEtwo:eqsecant}
Z_k^2G_k(te_0)\ge Z_k^2G_k(t_0e_0)\,e^{-(t-t_0)m_k(t_0)}\qquad(2\le t_0<t),
\end{equation}
whereas the lower half of (AF) asks for $Z_k^2G_k(te_0)\ge c_1w_\gamma(a_k
t)$, a power law, on $3\le t\le r_0/a_k$. The two are compared first as
functions of $t$ and then in spectral terms.

\begin{proposition}[exponential against power law]\label{rEtwo:compare}
\textup{(i)} If $\omega_k$ satisfies (AF), then for $3\le t$ with $(t+1)
a_k\le r_0$
\[
\Bigl|m_k(t)-8\log\frac{t+1}t\Bigr|\le\log\frac{C_1}{c_1}+\gamma\log
\frac{\ell(a_kt)}{\ell(a_k(t+1))} ,
\]
and the lattice mass $m_{\Phi,k}=\lim_tm_k(t)$ of Corollary~\ref{rEtwo:limit}
satisfies $\limsup_km_{\Phi,k}\le\log(C_1/c_1)$.

\textup{(ii)} Among all positive log-convex sequences with prescribed
values at $t_0$ and $t_0+1$, the geometric sequence $G(t_0)e^{-(t-t_0)
m(t_0)}$ is the smallest; so \eqref{rEtwo:eqsecant} is the best lower bound
at $t>t_0+1$ that follows from the values at $t_0$, $t_0+1$ and
log-convexity. If $G(t_0)$ and $G(t_0+1)$ are the values of the power law
$c\,t^{-8}$, this bound is $c\,t_0^{-8}(1+1/t_0)^{-8(t-t_0)}$, equal to
$c\,t^{-8}$ at $t=t_0+1$ and strictly smaller for every $t\ge t_0+2$.

\textup{(iii)} For every finite set $S\subset\{3,4,\dots\}$ and $c>0$ there
is a finite measure $\sigma$ on $[0,1]$ such that $G(t)=\int\lambda^{t-3}
d\sigma$ satisfies $\tfrac12c\,t^{-8}\le G(t)\le c\,t^{-8}$ for $t\in S$ and
$G(t)\le c\,\varepsilon^{t-3}$ for all $t\ge3$, for some $\varepsilon\in(0,1)$.
Hence the lower bounds $Z_k^2G_k(te_0)\ge c_1w_\gamma(a_kt)$ at the
separations $t\in S$, together with every inequality that the spectral
representation of Corollary~\ref{rEtwo:limit} imposes, do not imply the
lower bound at the separations $t\notin S$ with $t\le r_0/a_k$, once $a_k$
is small.
\end{proposition}

\begin{proof}
(i) is the computation of Proposition~\ref{rEtwo:degenerate}(iii) at the
separations $t,t+1$, where $\ell(a_kt)\ge\ell(a_k(t+1))$ since $\ell$ is
nonincreasing; the first alternative holds because the lower half of (AF)
gives $G_k(te_0)>0$ for $3\le t\le r_0/a_k$. For the mass, $m_{\Phi,k}\le
m_k(t)$ for every admissible $t$ by Corollary~\ref{rEtwo:limit}(iii);
taking $t+1=\lfloor r_0/a_k\rfloor$ and letting $k\to\infty$, $8\log((t+1)/
t)\to0$ and $\log(\ell(a_kt)/\ell(a_k(t+1)))\le\log(\ell(r_0-2a_k)/\ell(
r_0))\to0$.

(ii) A sequence $g(t)=\log G(t)$ is convex if and only if its increments
are nondecreasing, so $g(t)\ge g(t_0)+(t-t_0)(g(t_0+1)-g(t_0))$, with
equality for affine $g$, that is, geometric $G$. With $G(t_0)=ct_0^{-8}$
and $G(t_0+1)=c(t_0+1)^{-8}$ the geometric bound is $ct_0^{-8}(t_0/(t_0+1))
^{8(t-t_0)}$, and $(1+1/t_0)^{n}>1+n/t_0$ for $n=t-t_0\ge2$ is Bernoulli's
inequality, strict for $n\ge2$.

(iii) The power law is a moment sequence: $c(3+n)^{-8}=\int_0^\infty
e^{-(3+n)E}\,c\,E^7\,dE/7!$, so $c\,t^{-8}=\int\lambda^{t-3}d\sigma_0$ with
$\sigma_0$ the image of $cE^7e^{-3E}dE/7!$ under $\lambda=e^{-E}$. Let
$\sigma=\sigma_0|_{[\varepsilon',\varepsilon]}$ with $0<\varepsilon'<\varepsilon<1$.
Then $G(t)\le ct^{-8}$ for all $t$, $G(t)\le\varepsilon^{t-3}\sigma_0([0,1])
=c\,3^{-8}\varepsilon^{t-3}$, and $G(t)\ge ct^{-8}-\sigma_0([0,\varepsilon'))
-\sigma_0((\varepsilon,1])$ for every $t$. Since
$\sigma_0$ has no atoms, $\sigma_0([0,\varepsilon'))\to0$ as $\varepsilon'
\to0$ and $\sigma_0((\varepsilon,1])\to0$ as $\varepsilon\to1$; choose
$\varepsilon$ close to $1$ and then $\varepsilon'$ close to $0$ so that both
are at most $\tfrac14c(\max S)^{-8}$. Then $G(t)\ge\tfrac12ct^{-8}$ on $S$.
For the last statement take $c=2c_1a_k^{-8}$: then $G(t)\ge c_1a_k^{-8}
t^{-8}\ge c_1w_\gamma(a_kt)$ on $S$, $G$ has the spectral form of
Corollary~\ref{rEtwo:limit}(i), and $G(t)\le2c_1a_k^{-8}3^{-8}\varepsilon^{
t-3}<c_1w_\gamma(a_kt)$ for $t$ large with $a_kt\le r_0$, since
$\varepsilon^{t}t^{8}\ell(a_kt)^\gamma\to0$ and $\ell(a_kt)\le\ell(a_k)$.
\end{proof}

\paragraph{Spectral form of (AF).} Let $\sigma_k$ be the spectral measure
of Corollary~\ref{rEtwo:limit}(i) for $\omega_k$, and let $\tau_k$ be the
image of $Z_k^2\sigma_k$ under $\lambda\mapsto E=-\log\lambda\in[0,\infty]$,
the energy distribution, in lattice units of the time direction, of the
state $Z_k(\Phi_{2e_0}-c_k)\Omega'$ for the transfer matrix of the link
reflection. Then
\begin{equation}\label{rEtwo:eqlaplace}
Z_k^2G_k((3+n)e_0)=\int_{[0,\infty)}e^{-nE}\,d\tau_k(E)\qquad(n\ge1),
\end{equation}
and $\tau_k([0,\infty))\le\tau_k([0,\infty])=Z_k^2G_k(3e_0)$.

\begin{proposition}[spectral form of the two-point law]\label{rEtwo:spectral}
Let $\omega_k$ satisfy (AF), and put $A=\log(2C_1/c_1)$, $B=8+\gamma$.

\textup{(i)} \textup{(}Upper half.\textup{)} For $0<E\le1$,
\[
\tau_k([0,E])\le e\,C_1\,w_\gamma\bigl(a_k(3+\lfloor1/E\rfloor)\bigr)\le
e\,C_1\,a_k^{-8}E^{8} .
\]

\textup{(ii)} \textup{(}Lower half.\textup{)} For $4\le t\le r_0/a_k$ the
number $M_k(t)=\log\bigl(2C_1w_\gamma(3a_k)/(c_1w_\gamma(a_kt))\bigr)$
satisfies $\log2\le M_k(t)\le A+B\log(t/3)$ and
\[
\tau_k\Bigl(\Bigl[0,\frac{M_k(t)}{t-3}\Bigr]\Bigr)\ge\frac{c_1}2\,w_\gamma(
a_kt).
\]
Consequently, for $0<E\le1$ and $t_E=3+\lceil(2/E)(A+B\log(2B/E))\rceil$,
if $t_E\le r_0/a_k$ then
\begin{align*}
\tau_k([0,E])&\ge\frac{c_1}2\,w_\gamma(a_kt_E)\\
&\ge\frac{c_1}2\,a_k^{-8}\Bigl(\frac{E}{C_3\log(4/E)}\Bigr)^{8}\ell(a_k/E)
^{-\gamma},
\end{align*}
with $C_3=4+2A+2B\log(2B)+2B$; the condition $t_E\le r_0/a_k$ holds for
$E\ge C_4(a_k/r_0)\log(r_0/a_k)$ when $a_k/r_0$ is small, for a $C_4$
depending on $C_1/c_1$ and $\gamma$.

\textup{(iii)} \textup{(}Converse.\textup{)} If
\[
\tau_k([0,E])\ge C^{-1}a_k^{-8}\Bigl(\frac{E}{\log(4/E)}\Bigr)^8
\ell(a_k/E)^{-\gamma}
\]
for $E_k\le E\le1$, then for $11\le t\le3+8/E_k$
\[
Z_k^2G_k(te_0)\ge\frac{(8/e)^8}{C\,(1+\log16)^{\gamma}}\,\frac{w_\gamma(a_kt)}
{\bigl(\log((t-3)/2)\bigr)^8},
\]
the lower half of (AF) up to the factor $(\log((t-3)/2))^{-8}$; and if
$\tau_k([0,E])\le C'a_k^{-8}E^8$ for $0<E\le1$ and $\tau_k([0,\infty))\le
C'a_k^{-8}$, then $Z_k^2G_k(te_0)\le C''a_k^{-8}t^{-8}$ for $t\ge4$, the
upper half with $\gamma=0$.
\end{proposition}

\begin{proof}
(i) For $n\ge1$ and $E>0$,
\begin{align*}
\tau_k([0,E])&\le e^{nE}\!\int\! e^{-nE'}d\tau_k(E')
=e^{nE}Z_k^2G_k((3+n)e_0)\\
&\le e^{nE}C_1w_\gamma(a_k(3+n))
\end{align*}
by the upper half of (AF), valid at every $|z|_\infty\ge3$. With
$n=\lfloor1/E\rfloor\ge1$ one has $nE\le1$ and $3+n\ge1/E$, and
$w_\gamma(\rho)\le\rho^{-8}$.

(ii) By Corollary~\ref{rEtwo:limit}(i) and (AF),
\[
c_1w_\gamma(a_kt)\le Z_k^2G_k(te_0)\le Z_k^2G_k(3e_0)\le C_1w_\gamma(3a_k),
\]
so $M_k(t)\ge\log2$; and
\begin{align*}
\frac{w_\gamma(3a_k)}{w_\gamma(a_kt)}&=\Bigl(\frac t3\Bigr)^8
\Bigl(\frac{\ell(a_kt)}{\ell(3a_k)}\Bigr)^{\gamma},\\
\frac{\ell(3a_k)}{\ell(a_kt)}&=1+\frac{\log(t/3)}{\ell(a_kt)}\le1+\log(t/3),
\end{align*}
so $M_k(t)\le A+8\log(t/3)+\gamma\log(t/3)$. Split \eqref{rEtwo:eqlaplace} at
$E_0$ with $n=t-3\ge1$:
\begin{align*}
Z_k^2G_k(te_0)&\le\tau_k([0,E_0])+e^{-(t-3)E_0}\,\tau_k((E_0,\infty))\\
&\le\tau_k([0,E_0])+e^{-(t-3)E_0}\,C_1w_\gamma(3a_k).
\end{align*}
With $E_0=M_k(t)/(t-3)$ the last term is $\tfrac12c_1w_\gamma(a_kt)$ by the
definition of $M_k(t)$, and the lower half of (AF) gives the claim. For the
consequence: for $s\ge0$ and $0<\eta\le1$, $\log(1+s)\le\eta s+\log(1/
\eta)$, since the left side minus $\eta s$ is maximal at $s=1/\eta-1$ with
value $\log(1/\eta)-1+\eta$. With $\eta=E/(2B)$, $s=t-3$ and $t/3\le1+(t-3)$,
\[
M_k(t)\le A+B\log(1+(t-3))\le A+\tfrac12E(t-3)+B\log(2B/E),
\]
so $M_k(t)/(t-3)\le E$ as soon as $(t-3)E/2\ge A+B\log(2B/E)$, which holds
for $t=t_E$. If $t_E\le r_0/a_k$ the first bound applies at $t_E$ and
gives $\tau_k([0,E])\ge\tau_k([0,M_k(t_E)/(t_E-3)])\ge\tfrac12c_1w_\gamma(
a_kt_E)$. Finally $t_E\le4+(2/E)(A+B\log(2B)+B\log(1/E))\le C_3\log(4/E)/E$
for $E\le1$, using $1\le\log(4/E)$ and $\log(1/E)\le\log(4/E)$ for $E\le1$, and
$t_E\ge1/E$ gives $\ell(a_kt_E)\le\ell(a_k/E)$, $\ell$ being nonincreasing.
The condition $t_E\le r_0/a_k$ follows from $C_3\log(4/E)/E\le r_0/a_k$;
with $E=C_4y\log(1/y)$, $y=a_k/r_0$, this reads $C_3\log(4/(C_4y\log(1/y)))
\le C_4\log(1/y)$ and holds for $C_4=2C_3$ once $y$ is small, and
$\log(4/E)/E$ is decreasing in $E$.

(iii) For $t\ge11$ put $E=8/(t-3)\le1$; if $E\ge E_k$, that is $t\le3+8/
E_k$, then by \eqref{rEtwo:eqlaplace}
\begin{align*}
Z_k^2G_k(te_0)&\ge e^{-(t-3)E}\tau_k([0,E])\\
&\ge\frac{a_k^{-8}}{e^{8}C}\Bigl(\frac{8}{(t-3)\log((t-3)/2)}\Bigr)^8
\ell\Bigl(\frac{a_k(t-3)}8\Bigr)^{-\gamma},
\end{align*}
and $(t-3)/8\ge t/16$ for $t\ge6$ gives
\begin{align*}
\ell\bigl(a_k(t-3)/8\bigr)&\le\ell(a_kt/16)\le\ell(a_kt)+\log16\\
&\le(1+\log16)\,\ell(a_kt),
\end{align*}
while $(t-3)^{-8}\ge t^{-8}$.
For the upper half, integrate by parts in \eqref{rEtwo:eqlaplace} with
$n=t-3\ge1$:
\begin{align*}
Z_k^2G_k(te_0)&=\int_0^\infty ne^{-nE}\tau_k([0,E])\,dE\\
&\le C'a_k^{-8}\int_0^1ne^{-nE}E^8\,dE+e^{-n}C'a_k^{-8}\\
&\le C'a_k^{-8}(8!\,n^{-8}+e^{-n}),
\end{align*}
$e^{-n}\le8!\,n^{-8}$, and $(t-3)^{-8}\le4^8t^{-8}$ for $t\ge4$.
\end{proof}

\begin{remark}\label{rEtwo:spectralremark}
Proposition~\ref{rEtwo:spectral} says that (AF) along an axis is, up to
powers of logarithms, the statement that the energy distribution of the
local excitation $Z_k(\Phi_{2e_0}-c_k)\Omega'$ has weight of order
$a_k^{-8}E^8$ below every lattice energy $E$ between $(a_k/r_0)\log(r_0/a_k)$
and $1$: in physical units $E_{\mathrm{ph}}=E/a_k$, weight of order
$E_{\mathrm{ph}}^{\,8}$ below every energy between $\log(r_0/a_k)/r_0$ and
the cutoff $1/a_k$, the distribution $E_{\mathrm{ph}}^{\,7}dE_{\mathrm{ph}}$.
The lower half concerns the low-energy part at every intermediate scale;
the exponential bound \eqref{rEtwo:eqsecant} controls the Laplace transform
\eqref{rEtwo:eqlaplace} at a single $n$, hence, by
Proposition~\ref{rEtwo:compare}(ii),(iii), nothing of this kind.
\end{remark}

\paragraph{The two inputs of the lower half.} This subsection provides, for
the sequence $\omega_k$, two kinds of information: the positivity,
log-convexity and dichotomy of Theorem~\ref{rEtwo:axial} and
Corollary~\ref{rEtwo:limit}, valid at every $k$ with no uniformity in $k$,
and the explicit values of Theorem~\ref{rEtwo:strong}, valid for $\beta_k
\le\beta_2(N)$ only. A continuum sequence has $\beta_k\to\infty$ by
(AF$_\beta$) and leaves the second range; inside it,
Theorem~\ref{rEtwo:taylor} shows that the two-point function is
exponential in $t$, $G_k(te_0)\asymp\beta_k^{4t}$, in accordance with
Theorem~\ref{thm:limitEaf}(v).

\begin{corollary}[the two inputs]\label{rEtwo:step}
Let $\omega_k$ be torus limits at couplings $\beta_k\to\infty$ in the first
alternative of Corollary~\textup{\ref{rEtwo:limit}(ii)}, and consider

\textup{(a)} \textup{(}cutoff scale\textup{)} $\liminf_kZ_k^2G_k(3e_0)/
w_\gamma(3a_k)>0$, and

\textup{(b)} \textup{(}intermediate scales\textup{)} for all $k\ge k_0$
and $E_k\le E\le1$,
\[
\tau_k([0,E])\ge C^{-1}a_k^{-8}\Bigl(\frac{E}{\log(4/E)}\Bigr)^8
\ell(a_k/E)^{-\gamma}.
\]

(AF) implies \textup{(a)} and \textup{(b)} with $E_k=
C_4(a_k/r_0)\log(r_0/a_k)$. Conversely \textup{(a)} and \textup{(b)} imply
the lower half of (AF), with a constant factor on $3\le t\le10$ and up to the factor $(\log((t-3)/2))^{-8}$ on $11\le t
\le3+8/E_k$; with the $E_k$ above this is the range $a_kt\le8r_0/(C_4\log(
r_0/a_k))(1+o(1))$, shorter than $a_kt\le r_0$ by a logarithm. Neither
\textup{(a)} nor \textup{(b)} follows from the results of this subsection:
\textup{(a)} is a lower bound at three lattice units for the Wilson
measure at $\beta_k\to\infty$, where the strong-coupling expansion does
not reach, analyticity in $\beta$ gives no bound uniform in $k$, and
reflection positivity gives $G_k(3e_0)\ge0$ only; \textup{(b)} is a lower
bound on the low-energy spectral weight of the local excitation at every
scale between $r_0$ and the cutoff, which by
Proposition~\ref{rEtwo:compare}(iii) is not implied by \textup{(a)}, nor by
the lower half of (AF) at any finite set of separations, together with
every consequence of reflection positivity.
\end{corollary}

\begin{proof}
If (AF) holds, (a) is the case $t=3$ of its lower half and (b) is
Proposition~\ref{rEtwo:spectral}(ii) with $C=2C_3^8/c_1$. Conversely (b)
gives, by Proposition~\ref{rEtwo:spectral}(iii), the lower half for $11\le
t\le3+8/E_k$ up to the factor $(\log((t-3)/2))^{-8}$; (a) gives it at
$t=3$ since $w_\gamma(3a_k)\ge3^{-8}\ell(3a_k)^{-\gamma}a_k^{-8}$ is the
right side of (AF) there; and Corollary~\ref{rEtwo:limit}(i) gives it at
$4\le t\le10$ from the value at $t=11$, since $G_k(te_0)\ge G_k(11e_0)$
there and $w_\gamma(a_kt)\le(11/t)^8w_\gamma(11a_k)\le(11/4)^8w_\gamma(11a_k)$
for $4\le t\le11$, $\ell$ being nonincreasing. The negative statements are
Proposition~\ref{rEtwo:degenerate}(ii), Theorem~\ref{rEtwo:taylor},
Proposition~\ref{rEtwo:rp} and Proposition~\ref{rEtwo:compare}(iii).
\end{proof}

\begin{remark}\label{rEtwo:balaban}
Statement (b) at the scale $E$ concerns, in the renormalisation-group
description of Ba{\l}aban~\cite{Bal87,Bal88,Bal89b}, the effective theory
after the block-spin steps that bring the lattice spacing from $a_k$ to
about $a_k/E$, with the two insertions of $\Phi$ made at the cutoff scale
and carried through every step. As recalled in
Remarks~\ref{rem:limitEafstatus} and~\ref{rE:status}, the bounds
of~\cite{Bal89b} concern the partition function of a box of fixed physical
size, with the large-field regions localised and exponentiated, and give
no lower bound on a connected correlation. What Corollary~\ref{rEtwo:step}
adds to those remarks is the form of the missing bound: a lower bound on
the low-energy weight of the two-source generating function at every
intermediate scale, with the large-field contribution dominated by the
small-field connected part of order $a_k^{-8}E^8$ at that scale, and the
same at the cutoff scale for (a). Positivity, analyticity and the
strong-coupling expansion give the lower bound at every scale from the
values at two consecutive separations, but the values at two separations,
at any fixed scale, do not determine the low-energy weight at the others.
\end{remark}

\begin{remark}[status]\label{rEtwo:status2}
Proved here, for every $\beta>0$, every even torus $L\ge6$ and every
compact $G\subset\U(N)$ with $\rho$ irreducible: the axial connected
correlation of $\Phi$ is nonnegative, log-convex and monotone, and either
strictly positive at all axial separations or zero at all of them, the
latter only on a closed discrete set $\mathcal D_L$ of couplings that
misses $(0,\beta_1(N)]$; for a single spatial plaquette it is strictly
positive at every $\beta$; the exponential lower bound
\eqref{rEtwo:eqsecant} holds with the exact rate $m_\Phi$ in torus limits;
its constants are $c(\beta)=108((1+\delta)/2)^9\beta^8$ and $m(\beta)=4
\log(2/((1+\delta)\beta))$ to leading order at strong coupling; reflection
positivity propagates positivity but does not create it. Not proved: that
$\mathcal D_L$ is empty for every $L$; positivity at non-axial
separations; any bound uniform in the lattice spacing along a continuum
sequence; either half of (AF). Nothing here bears on the existence of a
continuum limit or on a mass gap; Corollary~\ref{rEtwo:step} states the two
inputs that an argument for the lower half of (AF) along an axis would
have to supply.
\end{remark}
